\paragraph{AQ-136-HIGHER-COMPLETION: the positive partial axioms do not ensure relative completion.}
The partial source of AQ-136-PARTIAL-SOURCE is not claimed to
admit an arbitrary further relative repair. The following example
shows that such a repair theorem would be false, even without
requiring the completed source to have simplex zero exactly $S$.
It does not refute the original common-source conjecture.

Fix a rational prime $p$, and work over $k=\mathbb Z$.
There is a simplicial strict PD differential $S$-algebra $A$
and two maps $f_0,f_1:A\to B$ to the standard admissible
theory, with the following properties:
\[
 \begin{gathered}
 A_0=S,\qquad A_n^{0,*}=\mathbb Z[0],\\
 H^0(A_n^{w,*})=S^w,\qquad H^1(A_n^{w,*})=0
                                      \quad\text{for all }n,w,\\
 H^*(A_{\mathbb Q,n})=S_{\mathbb Q},\qquad
 H_r(N A_*^{w,q})=0\quad(w>q,\ r\ge0).
 \end{gathered}                                            \tag{136.HC1}
\]
Nevertheless there are no admissible $R$ and strict maps
\[
            A\xrightarrow{g}R
                    \mathrel{\substack{\xrightarrow{j_0}\\[-3pt]
                                      \xrightarrow[j_1]{}}}B,
                 \qquad j_i g=f_i.                       \tag{136.HC2}
\]
Both targets are the same object $B$, so they certainly have
a common admissible source. The obstruction is precisely to
extending the specified maps from this specified partial source.

\begin{proof}
Let $L_n=\mathbb Z^{n+1}$ be the standard coordinate simplicial
module, with coefficient inclusion $\mathbb Z\to L_n$ given
by the sum of the coordinate generators. Put
\[
             Q=L/\mathbb Z,\qquad NQ=\mathbb Z[1].
\]
These are the modules already used in the construction of $B$.
Let $V$ and $U$ be copies of $Q$, placed in form degrees one
and two respectively, and give both weight one. Define
\[
 D=\Lambda_{\mathbb Z}(V)\otimes_{\mathbb Z}\Gamma_{\mathbb Z}(U),
        \qquad d v=p u,\quad d u=0,\qquad
 A=B\otimes_{\mathbb Z}D.                                 \tag{136.HC3}
\]
The correspondence $v\mapsto u$ is the identity between the
two copies of $Q$. Thus the differential is simplicial and
respects every linear relation. Extend it by the strict
graded derivation and special PD rules. Its square is zero;
the odd exterior relations are stable: $d(v^2)=p(uv-vu)=0$,
and the same derivation rule preserves $vv'+v'v=0$.
The usual free graded PD algebra structure supplies divided
powers on positive even elements. The coefficient map to $A$
is its map to the $B$ factor. This verifies the actual strict
PD structure, not merely its underlying bicomplex.

In simplicial degree $n$, the module $Q_n$ is free of rank
$n$. Hence $D_n^0=\mathbb Z$, and
\[
     d:D_n^1=\bigoplus_{i=1}^n\mathbb Z v_i
             \longrightarrow D_n^2
\]
sends the basis $v_i$ to the linearly independent elements
$p u_i$. It is injective. Consequently
\[
                       H^0(D_n)=\mathbb Z,
                              \qquad H^1(D_n)=0.           \tag{136.HC4}
\]
The standard integral PD de Rham contraction of $B_n$ onto
$S$ is $S$-linear and weight-preserving: after solving for
the zeroth coordinate, it is the tensor product of the
one-variable contractions
$x^{[m-1]}dx\mapsto x^{[m]}$ for $m\ge1$.
Tensoring that homotopy with $D_n$ identifies the cohomology
of $A_n$ with that of $S\otimes D_n$.
The module $S$ is free over $\mathbb Z$, so (136.HC4) proves
the asserted integral $H^0$ and $H^1$ in every weight.
Also $Q_0=0$, so $D_0=\mathbb Z$ and $A_0=B_0=S$.
Positive form generators have positive weight, which proves
the weight-zero assertion.

Over $\mathbb Q$ the algebra on each pair $v_i,u_i$ contracts
onto its constants. Indeed
\[
             d(v_i u_i^{[m]})=p(m+1)u_i^{[m+1]},
\]
and $p(m+1)$ is invertible. The tensor product of these
contractions proves the rational vertical assertion.
In particular all positive integral vertical cohomology of
$A_n$ is integer-torsion, as in the partial-source theorem.

For the horizontal assertion, a monomial of form degree $d$
in $D$ has weight $e$ satisfying
\[
                          \lceil d/2\rceil\le e\le d.
                                                               \tag{136.HC5}
\]
There are only finitely many possible $e$ for fixed $d$.
Every summand of $A^{w,q}$ therefore has the form
\[
                 B^{w-e,q-d}\otimes D^{e,d}.
\]
If $w>q$, then $w-e>q-d$. The standard row contraction
of $B^{a,b}$ for $a>b$ applies: its positive coordinate
PD degree is $a-b>0$, and the linear contraction of $L$
induces the contraction of that PD power. Tensoring in
the simplicial direction preserves this acyclicity.
Equivalently, the normalized Eilenberg--Zilber comparison
identifies the diagonal tensor with the tensor of normalized
complexes, whose first factor is contractible. All modules
here are degreewise free; the same flat comparison used in
AQ-136-TENSOR applies (Stacks Project, Remark~21.39.11,
Tag~08QD). Summing the finitely many summands proves the
last assertion of (136.HC1), including horizontal $H_0$.

The standard theory contains the canonical copy of $Q$ as
its closed one-form generators. Define $f_0$ and $f_1$ to
be the identity on $B$, and on $D$ by
\[
       f_0(v)=f_0(u)=0,\qquad f_1(v)=e,
                                  \qquad f_1(u)=0.        \tag{136.HC6}
\]
Here $e$ denotes that canonical closed coordinate
one-form, not a globally chosen primitive on a simplicial
circle. More concretely its normalized generator in
simplicial degree one is $dx$, with coordinate $x$ having
endpoints $0,T$. The maps of copies of $Q$ are simplicial;
$d(dx)=0$ verifies compatibility with $dv=pu$.
The free exterior/PD universal property proves that
(136.HC6) defines strict maps in the stated category.

Let $X=\Delta[1]/\partial\Delta[1]$ be the simplicial circle.
The normalized generators $v$ and $u$ of the copies of $Q$
have both faces zero, so they define elements of $A(X)$,
still denoted $v,u$, satisfying $dv=pu$.
For every $w\ge1$ put
\[
             v_w=T^{[w-1]}v,\qquad
             u_w=T^{[w-1]}u,\qquad dv_w=p u_w.             \tag{136.HC7}
\]
On the edge, the difference $f_1(v_w)-f_0(v_w)$ is the
closed one-form $T^{[w-1]}dx$.
Its class in $H^1(B^{w,*}(X))=S^w$ is
\[
                  T^{[w-1]}T=wT^{[w]}.                    \tag{136.HC8}
\]
For clarity, this normalization can be checked without a
comparison convention: the integral PD antiderivative on
the edge sends $x^{[m-1]}dx$ to $x^{[m]}$.
The difference between its endpoint values identifies
$H^1(B^{w,*}(X))$ with $S^w$; exact circle forms have zero
endpoint difference, and $x^{[w-1]}dx$ has difference
$T^{[w]}$, so the map is an isomorphism. Formula (136.HC8)
follows from the divided-power multiplication rule.

Suppose (136.HC2) existed. The restriction $b:B\to R$ of
$g$ satisfies $j_0b=j_1b=1_B$. The natural comparison theorem
AQ-136-COMPARISON, applied through form degree two, gives,
for every sufficiently large weight $w$,
\[
 H^1(R^{w,*}(X))\simeq S^w,\qquad
 H^2(R^{w,*}(X))=0,\qquad
 H^1(b)\text{ is an isomorphism}.                         \tag{136.HC9}
\]
Thus $H^1(j_0)=H^1(j_1)$, both being the inverse of $H^1(b)$.
Choose such a $w$ not divisible by $p$. Since $u_w$ is
closed, the middle assertion of (136.HC9) provides
$a\in R^{w,1}(X)$ with $da=g(u_w)$. Then
\[
                           c=pa-g(v_w)
\]
is closed. Its two images have the same cohomology class.
The elements $j_i(a)$ are closed because $f_i(u_w)=0$.
Using (136.HC8), the equality of the two classes says
\[
        p\bigl([j_1(a)]-[j_0(a)]\bigr)=wT^{[w]}
                                      \quad\text{in }S^w.
                                                               \tag{136.HC10}
\]
But $S^w=\mathbb ZT^{[w]}$, and $p\nmid w$, which is a
contradiction. Arbitrarily large such weights exist, so
no choice of admissibility cutoff avoids it.
This proves the exact relative-completion obstruction.
\end{proof}

\paragraph{Consequence for the next construction step.}
The hypotheses of the positive partial-source theorem alone
cannot justify a functorial repair killing all higher
vertical torsion while retaining arbitrary prescribed target
maps. Such a construction would have to maintain an additional
compatibility invariant that excludes (136.HC3)--(136.HC6),
or replace the chosen partial source or its maps. No negative
answer to either original integral question is deduced here.
